\chapter{Séries temporelles}\label{ch:iv1:time-series}

\begin{ivtrains}
Les outils de travail de la modélisation des rendements et de la volatilité~: stationnarité, AR et MA, GARCH, estimateurs de volatilité par l'étendue et à haute fréquence, racines unitaires et cointégration, et manière de comparer les prévisions. Les questions demandent un nombre et la raison pour laquelle ce nombre induit en erreur.
\end{ivtrains}

\begin{question}[$\star$ \iqroles{QR, QT, DS, Risque} \iqfirm{toutes} \iqlevel{junior}]\label{iq:iv1:time-series:1}
Lesquelles de ces suites sont faiblement stationnaires~? (a) Un bruit blanc. (b) Une marche aléatoire. (c) Un AR(1) de coefficient $0.9$, démarré dans sa loi stationnaire. (d) Une tendance linéaire plus un bruit. (e) La différence première d'une marche aléatoire. (f) $x_t = A\cos(t) + B\sin(t)$, où $A$ et $B$ sont des normales centrées réduites indépendantes, tirées une seule fois.
\end{question}

\begin{question}[$\star$ \iqroles{QR, QT, DS, Risque} \iqfirm{toutes} \iqlevel{junior}]\label{iq:iv1:time-series:2}
Un AR(1) vérifie $x_t = 0.7x_{t-1} + \varepsilon_t$, avec une variance du bruit égale à 1. Quelles sont la variance et l'autocorrélation au retard~3~? Quelle est la demi-vie d'un choc~? Comment le $R^2$ de la prévision à $h$ pas dépend-il de $h$, et quelle est sa valeur pour $h = 1$ et $h = 2$~?
\end{question}

\begin{question}[$\star$ \iqroles{QR, QT, DS, Risque} \iqfirm{toutes} \iqlevel{junior}]\label{iq:iv1:time-series:3}
Un MA(1) vérifie $x_t = \varepsilon_t + 0.5\varepsilon_{t-1}$, avec une variance du bruit égale à 1. Donnez la variance et les autocorrélations aux retards 1 et 2. Quelle est la plus grande autocorrélation possible au retard 1 pour un MA(1)~? Quelle autre valeur du coefficient donne les mêmes autocorrélations~?
\end{question}

\begin{question}[$\star$ \iqroles{QR, QT, DS, Risque} \iqfirm{toutes} \iqlevel{junior}]\label{iq:iv1:time-series:4}
Une variance exponentiellement pondérée utilise un facteur de décroissance de $0.94$ sur les rendements quotidiens au carré. Quel poids reçoit le dernier jour~? Quelle est la demi-vie des poids, et quel est le nombre effectif de jours~? Un mouvement important, égal à 5 fois la volatilité quotidienne courante, survient. De quel facteur l'estimation de volatilité saute-t-elle~?
\end{question}

\begin{question}[$\star\star$ \iqroles{QR, QT, DS} \iqfirm{toutes} \iqlevel{junior, confirmé}]\label{iq:iv1:time-series:5}
Considérez $x_t = 1.2x_{t-1} - 0.5x_{t-2} + \varepsilon_t$. Ce processus est-il stationnaire~? Trouvez les racines du polynôme caractéristique, le module et la période d'un éventuel cycle, ainsi que les autocorrélations aux retards 1 et 2. Quel serait l'effet d'un choc au fil du temps~?
\end{question}

\begin{question}[$\star\star$ \iqroles{QR, QT, DS} \iqfirm{toutes} \iqlevel{junior}]\label{iq:iv1:time-series:6}
Une série de 500 points présente des autocorrélations empiriques $0.60$, $0.35$, $0.22$, $0.12$ et des autocorrélations partielles empiriques $0.60$, $-0.02$, $0.01$, $0.03$. Quel modèle proposez-vous, et quelles bornes de confiance utilisez-vous pour juger de la significativité~? Quelle forme présenterait un MA(1), et que donnerait un ARMA(1,1)~?
\end{question}

\begin{question}[$\star\star$ \iqroles{QR, QT, DS} \iqfirm{hedge fund, banque} \iqlevel{junior, confirmé}]\label{iq:iv1:time-series:7}
Les rendements suivent un AR(1) d'autocorrélation $0.1$. Quel est le ratio de variance $\operatorname{Var}(r_t + \dots + r_{t-q+1})/(q\operatorname{Var}r_t)$ pour $q = 2$ et $q = 5$~? Quelle est la limite pour $q$ grand~? De combien l'annualisation usuelle en racine carrée du temps sous-estime-t-elle la volatilité~?
\end{question}

\begin{question}[$\star\star$ \iqroles{QR, QT, DS, Risque} \iqfirm{hedge fund, banque} \iqlevel{junior, confirmé}]\label{iq:iv1:time-series:8}
Un GARCH(1,1) vérifie $\sigma_t^2 = \omega + \alpha r_{t-1}^2 + \beta\sigma_{t-1}^2$, avec $\alpha = 0.1$, $\beta = 0.85$ et $\omega = 2 \times 10^{-6}$. Trouvez la variance de long terme et la volatilité annuelle, sur 252 jours. En combien de temps un choc de volatilité est-il divisé par deux~? Si la variance d'aujourd'hui vaut quatre fois la variance de long terme, quelle est la prévision à 10 jours~?
\end{question}

\begin{question}[$\star\star$ \iqroles{QR, QT, DS, Risque} \iqfirm{hedge fund, teneur de marché} \iqlevel{junior, confirmé}]\label{iq:iv1:time-series:9}
Le cours le plus haut et le cours le plus bas d'une journée sont dans le rapport $H/L = 1.02$. Que donne l'estimateur de Parkinson $(\ln(H/L))^2/(4\ln 2)$ comme volatilité quotidienne~? De combien est-il plus efficace qu'un seul rendement entre deux clôtures successives, au carré~? Donnez les deux principales façons dont il se trompe sur des données réelles.
\end{question}

\begin{question}[$\star\star$ \iqroles{QR, QT} \iqfirm{teneur de marché, hedge fund} \iqlevel{junior, confirmé}]\label{iq:iv1:time-series:10}
Le logarithme du prix observé est égal au prix efficient augmenté d'un bruit indépendant d'écart type $0.05~\%$, et le prix efficient a une volatilité quotidienne de 1~\%. Quelle est l'espérance de la variance réalisée calculée à partir de $n$ rendements intrajournaliers de même longueur~? Quel est le biais relatif avec 78 rendements par jour, et avec 390~? Comment le bruit se manifeste-t-il dans l'autocovariance des rendements~?
\end{question}

\begin{question}[$\star\star$ \iqroles{QR, QT, DS} \iqfirm{hedge fund, banque} \iqlevel{junior, confirmé}]\label{iq:iv1:time-series:11}
Vous estimez $x_t = \rho x_{t-1} + \varepsilon_t$ alors qu'en vérité $\rho = 1$. Pourquoi la statistique $t$ usuelle sur $\hat{\rho} - 1$ n'utilise-t-elle pas la table normale~? Quelles sont les valeurs critiques à 5~\% sans constante et avec une constante~? Avec une constante et 250 observations, quelle valeur de $\hat{\rho}$ attendez-vous, à peu près~? Montrez ces résultats par simulation.
\end{question}

\begin{question}[$\star\star$ \iqroles{QR, QT, DS} \iqfirm{hedge fund} \iqlevel{junior, confirmé}]\label{iq:iv1:time-series:12}
Deux prévisions sont comparées sur 400 périodes. La différence de leurs erreurs quadratiques a pour moyenne $0.02$ et pour écart type $0.2$. Trouvez la statistique de Diebold--Mariano si ces différences sont indépendantes dans le temps, puis si elles suivent un AR(1) de coefficient $0.2$. Qu'est-ce qui change si l'un des modèles a été retenu comme le meilleur parmi de nombreux~?
\end{question}

\begin{question}[$\star\star$ \iqroles{QR, QT} \iqfirm{hedge fund} \iqlevel{junior, confirmé}]\label{iq:iv1:time-series:13}
$x_t$ est une marche aléatoire, et $y_t = 2x_t + u_t$, où $u_t$ est un AR(1) de coefficient $0.9$. Quels sont les ordres d'intégration de $x$, $y$ et $u$~? Quelle est la demi-vie du spread, et à quelle vitesse la pente de la régression converge-t-elle~? Décrivez les deux étapes de la procédure d'Engle--Granger, et expliquez pourquoi la seconde étape ne peut pas utiliser la table usuelle. En quoi cela diffère-t-il de la régression de deux marches aléatoires indépendantes~?
\end{question}

\begin{question}[$\star\star\star$ \iqroles{QR, QT, DS} \iqfirm{hedge fund, banque} \iqlevel{confirmé}]\label{iq:iv1:time-series:14}
Pour le GARCH(1,1) à innovations gaussiennes, $\alpha = 0.1$ et $\beta = 0.85$, le kurtosis vaut $3(1 - (\alpha + \beta)^2)/(1 - (\alpha + \beta)^2 - 2\alpha^2)$. Dérivez ce résultat en calculant le moment d'ordre 4. Évaluez-le. Quelle est la condition pour qu'il soit fini~? Les rendements eux-mêmes sont non corrélés. Quelle est l'autocorrélation au retard 1 des rendements au carré, donnée par $\alpha(1 - \alpha\beta - \beta^2)/(1 - 2\alpha\beta - \beta^2)$, et à quelle vitesse décroît-elle~?
\end{question}

\begin{question}[$\star\star\star$ \iqroles{QR, QT, DS} \iqfirm{hedge fund, banque} \iqlevel{confirmé}]\label{iq:iv1:time-series:15}
Une prévision parfaite de la variance conditionnelle $\sigma_t^2$ est confrontée au rendement au carré $r_t^2$. Montrez que le $R^2$ de cette régression est au plus égal à $(\kappa/3 - 1)/(\kappa - 1)$, où $\kappa$ est le kurtosis du rendement et où les innovations sont gaussiennes. Évaluez ce résultat pour le GARCH(1,1) de kurtosis $3.77$. Qu'en dit-il pour juger des modèles de volatilité~? Comment fixez-vous la comparaison~?
\end{question}

\begin{question}[$\star\star\star$ \iqroles{QR, QT} \iqfirm{hedge fund} \iqlevel{confirmé}]\label{iq:iv1:time-series:16}
Un rendement vérifie $r_t = \mu_t + \varepsilon_t$, où le rendement espéré $\mu_t$ est un AR(1) lent de coefficient $\phi = 0.9$ et où le bruit est un bruit blanc indépendant de lui. Le signal représente $R^2 = 1~\%$ de la variance totale. Quelle est l'autocorrélation au retard 1 des rendements, et au retard $k$~? Quel modèle ARMA suit $r$~? Pourquoi les autocorrélations des rendements peuvent-elles paraître nulles alors que la prévisibilité est réelle~?
\end{question}

\begin{question}[$\star\star\star$ \iqroles{QR, QT, DS, Risque} \iqfirm{hedge fund, banque} \iqlevel{confirmé}]\label{iq:iv1:time-series:17}
Les rendements sont gaussiens conditionnellement à un état, avec une volatilité annualisée de 10~\% dans un état calme et de 30~\% dans un état de stress. L'état suit une chaîne de Markov qui reste dans l'état calme avec une probabilité de $0.95$, et dans l'état de stress avec une probabilité de $0.9$. Trouvez les probabilités stationnaires, la volatilité inconditionnelle et le kurtosis. Comment décroît l'autocorrélation des rendements au carré~?
\end{question}

\begin{question}[$\star\star\star$ \iqroles{QR, QT, DS} \iqfirm{hedge fund} \iqlevel{junior, confirmé}]\label{iq:iv1:time-series:18}
Vous ajustez un AR(1) par moindres carrés avec une constante sur 100 points d'un spread à retour à la moyenne, dont le vrai coefficient vaut $0.95$. Utilisez la formule du biais $\mathbb{E}\hat{\phi} - \phi \approx -(1 + 3\phi)/T$ pour calculer l'estimation espérée. Quel effet cela a-t-il sur la demi-vie que vous publiez~? Pourquoi ce biais apparaît-il, et comment le corrigeriez-vous~?
\end{question}

\begin{question}[$\star\star\star\star$ \iqroles{QR, QT} \iqfirm{hedge fund} \iqlevel{confirmé}]\label{iq:iv1:time-series:19}
Sous une racine unitaire sans constante, dérivez les lois limites de $T(\hat{\rho} - 1)$ et de la statistique $t$, en tant que fonctionnelles d'un mouvement brownien sur $[0, 1]$. Déterminez la loi du numérateur, ainsi que la probabilité qu'il soit négatif. Vérifiez par simulation la valeur critique à 5~\% de la statistique $t$ et la moyenne de $T(\hat{\rho} - 1)$.
\end{question}

\begin{question}[$\star\star\star\star$ \iqroles{QR, QT} \iqfirm{teneur de marché, hedge fund} \iqlevel{confirmé}]\label{iq:iv1:time-series:20}
Le prix efficient a une variance quotidienne $\sigma^2$, et le bruit a une variance $\eta^2$, comme dans un modèle standard de bruit de microstructure. Écrivez l'erreur quadratique moyenne de la variance réalisée en fonction de $n$, en gardant les termes principaux, et trouvez le $n$ optimal, pour $\sigma = 1~\%$ et $\eta = 0.05~\%$. Montrez ensuite comment la moyenne sur des sous-échantillons, combinée à la soustraction d'un terme de biais, donne un estimateur qui utilise toutes les données et n'a pas de biais du premier ordre, puis testez-le par simulation.
\end{question}
